% Generated draft: requires mathematical and source review.
\documentclass[11pt]{article}
\usepackage[margin=25mm]{geometry}
\usepackage{fontspec}
\setmainfont{Latin Modern Roman}
\usepackage{amsmath,amssymb,mathtools,unicode-math}
\setmathfont{Latin Modern Math}
\usepackage{enumitem,needspace,xurl,hyperref}
\setlength{\parindent}{0pt}
\setlength{\parskip}{5pt}
\setlength{\emergencystretch}{3em}
\newcommand{\N}{\mathbb{N}}
\newcommand{\Z}{\mathbb{Z}}
\newcommand{\Q}{\mathbb{Q}}
\newcommand{\R}{\mathbb{R}}
\newcommand{\C}{\mathbb{C}}
\newcommand{\F}{\mathbb{F}}
\newcommand{\A}{\mathbb{A}}
\newcommand{\PP}{\mathbb{P}}
\newcommand{\E}{\mathbb{E}}
\newcommand{\eps}{\varepsilon}
\newcommand{\dd}{\,\mathrm{d}}
\begin{document}
\section*{UnsolvedMath ID: 30000521}
\textbf{Source problem code:} catalogue entry
\subsection{Definitions and mathematical statement}
Let $G$ be finite and $k$ a field of characteristic different from $2$. A $G$-Galois algebra $L$ is a finite \'etale $k$-algebra with $G$-action which becomes the regular $G$-permutation algebra after base change to a separable closure; in particular $\dim_kL=|G|$. Its trace $G$-form is
\[
 q_L(x)=\operatorname{Tr}_{L/k}(x^2),
 \qquad b_L(x,y)=\operatorname{Tr}_{L/k}(xy),
\]
where $b_L$ is the associated half-polarization and $G$ acts by algebra automorphisms. The form is $G$-invariant.

For two such algebras $L,L'$, the first question is to characterize when there is a $G$-equivariant $k$-linear isomorphism $T:L\to L'$ satisfying
\[
 q_{L'}(T x)=q_L(x)\quad\forall x\in L.
\]
It is not required to preserve multiplication, and forgetting $G$ would give a weaker ordinary quadratic-form question.

The second asks, for the given $L,L'$, which quadratic forms $\varphi$ over $k$ make
\[
 \varphi\otimes q_L\simeq_G\varphi\otimes q_{L'}.
\]
On the tensor product the $\varphi$-space has trivial $G$-action. In the usual nonsingular finite-dimensional setting the tensor form is defined by the tensor product of the associated symmetric bilinear forms. The excerpt does not specify whether degenerate or zero forms are included. The desired conclusion is an actual equivariant isometry, not merely an equality after forgetting the action or an unstated stabilization/cancellation convention.

The supplied literature triage checked 21 August 2026 reports a local--global isomorphism theorem for $G$-trace forms over global fields by Bayer--Fluckiger, Parimala and Serre, while leaving the unrestricted field/characterization question open. A theorem over global fields does not classify all fields of characteristic different from $2$, and the tensor-multiplier classification is a separate request. The exact full theorem hypotheses and its application to the weaker question are not reproduced, so no broader result is inferred.
\subsection{Short English statement}
When are the trace forms of two $G$-Galois algebras isometric with the $G$-action respected, and which quadratic-form tensor factors make them isometric? The supplied assessment reports a global-field Hasse principle, without answering the unrestricted field and multiplier classification questions.
\subsection{Sources}
\begin{itemize}
\item Local UnsolvedMath registry, permanent ID 30000521, exact source code catalogue entry; \emph{Isomorphism of Trace Forms for $G$-Galois Algebras}. Proposed year 2006, status partially\_solved; corrected\_verified extraction retains both questions and characteristic not $2$.
\item Supplied report citation: \emph{Quadratic Forms and Linear Algebraic Groups (2007)}, Oberwolfach Reports, DOI 10.4171/owr/2006/29, pp. 1743--1794, \url{https://doi.org/10.4171/owr/2006/29}. Title year differs from DOI/proposed year.
\item Supplied literature triage checked 2026-08-21: ``original Oberwolfach report'', \url{https://doi.org/10.4171/owr/2006/29}; ``Hasse principle for $G$-trace forms'', \url{https://arxiv.org/abs/1210.5765}, attributed in the assessment to Bayer--Fluckiger, Parimala and Serre. This dated OWR-derived curation requests independent verification of time-sensitive status and cited claims.
\end{itemize}
\end{document}
